\section{Interpretation}

The current experiments show that the direct diffusion model provides the
strongest unconditional generative baseline in terms of overall sample quality
and agreement with the validation distribution. The wavelet-based formulation
separates into two regimes: its conditioned detail stages are accurate when the
coarse state is supplied, while the fully generated cascade is presently
limited by coarse-stage error propagation.

The comparison therefore does not simply oppose a good model to a bad one. It
highlights a genuine methodological trade-off: direct diffusion currently gives
better unconditional samples, whereas wavelet diffusion offers a structured
coarse-to-fine representation whose potential for fast sampling depends on a
more accurate unconditional coarse model.

% This evaluation chapter showed the strengths and limitations of the implemented pipelines. The final chapter summarizes the main findings and discusses perspectives for improving the wavelet-based approach.
